%% =======================================================
%% Actividad 2:
%%
%%
%% =======================================================

%% ---- CONFIGURACIÓN ------------------------

\documentclass[11pt,letterpaper]{article}

\usepackage[indice]{tecnm}
\usepackage[T1]{fontenc}
\usepackage{xcolor}
\usepackage[most]{tcolorbox} % Carga tcolorbox con soporte para listings y saltos de página
\tcbuselibrary{listings, skins, breakable}

% Paleta de colores (estilo GitHub Light / VS Code)
\definecolor{codebg}{HTML}{F8F9FA}       % Fondo del código
\definecolor{codehead}{HTML}{E9ECEF}     % Fondo de la barra de título
\definecolor{codeborder}{HTML}{CED4DA}   % Borde de la tarjeta
\definecolor{codetitle}{HTML}{212529}    % Texto del título
\definecolor{codekw}{HTML}{0550AE}       % Palabras clave (import, def, for...)
\definecolor{codefunc}{HTML}{8250DF}     % Funciones y métodos (describe, groupby...)
\definecolor{codestr}{HTML}{0A3069}      % Strings
\definecolor{codecomm}{HTML}{6E7781}     % Comentarios
\definecolor{codenum}{HTML}{8C959F}      % Números de línea
\definecolor{codeid}{HTML}{24292F}       % Texto general

% Configuración base de listings
\lstdefinestyle{pythonestilo}{
  language=Python,
  basicstyle=\ttfamily\footnotesize\color{codeid},
  keywordstyle=\color{codekw}\bfseries,
  keywordstyle=[2]\color{codefunc}, % Estilo secundario para métodos de datos
  stringstyle=\color{teal!80!black},
  commentstyle=\color{codecomm}\itshape,
  morekeywords={as, True, False, None, with, yield, lambda},
  morekeywords=[2]{read_csv, head, tail, info, describe, groupby, mean, value_counts, to_csv, print, len, range},
  numbers=left,
  numberstyle=\tiny\ttfamily\color{codenum},
  stepnumber=1,
  numbersep=10pt,
  showstringspaces=false,
  morestring=[b]{"""},
  tabsize=4,
  breaklines=true,
  breakatwhitespace=true,
  columns=flexible,
  keepspaces=true,
  literate=
    {á}{{\'a}}1 {é}{{\'e}}1 {í}{{\'i}}1 {ó}{{\'o}}1 {ú}{{\'u}}1
    {Á}{{\'A}}1 {É}{{\'E}}1 {Í}{{\'I}}1 {Ó}{{\'O}}1 {Ú}{{\'U}}1
    {ñ}{{\~n}}1 {Ñ}{{\~N}}1 {¿}{{?`}}1 {¡}{{!`}}1
}

% Definición del entorno personalizado \begin{codigo}{Título}{etiqueta}
% \newtcblisting[auto counter, number within=section]{codigo}[2][]{
\newtcblisting[auto counter]{codigo}[2][]{
  listing only,
  breakable,                    % Permite que el bloque se divida si pasa de página
  enhanced,
  colback=codebg,
  colframe=codeborder,
  colbacktitle=codehead,
  coltitle=codetitle,
  fonttitle=\bfseries\small\sffamily,
  boxrule=0.6pt,
  arc=4pt,                      % Bordes redondeados
  left=20pt,                    % Espacio izquierdo para los números de línea
  right=8pt,
  top=6pt,
  bottom=6pt,
  toptitle=4pt,
  bottomtitle=4pt,
  before skip=12pt,
  after skip=12pt,
  listing options={style=pythonestilo},
  title={Código \thetcbcounter: #2},
  #1
}

% Paleta para salida de consola (estilo terminal limpia)
\definecolor{termbg}{HTML}{212529}      % Fondo oscuro tipo terminal
\definecolor{termhead}{HTML}{343A40}    % Barra superior
\definecolor{termtext}{HTML}{F8F9FA}    % Texto claro de alta legibilidad
\definecolor{termborder}{HTML}{495057}  % Borde

% Entorno \begin{consola}[label=out:iris]{Título}
\newtcblisting[auto counter]{consola}[2][]{
  listing only,
  breakable,
  colback=termbg,
  colframe=termborder,
  colbacktitle=termhead,
  coltitle=termtext,
  colupper=termtext,
  fonttitle=\bfseries\small\sffamily,
  boxrule=0.6pt,
  arc=4pt,
  left=8pt,
  right=8pt,
  top=6pt,
  bottom=6pt,
  toptitle=4pt,
  bottomtitle=4pt,
  before skip=12pt,
  after skip=12pt,
  listing options={
    basicstyle=\ttfamily\scriptsize\color{termtext}, % \scriptsize evita que las tablas de Pandas se partan de línea
    breaklines=true,
    columns=fixed,        % Mantiene alineadas las columnas de Pandas
    keepspaces=true,
    showstringspaces=false,
    literate=
      {á}{{\'a}}1 {é}{{\'e}}1 {í}{{\'i}}1 {ó}{{\'o}}1 {ú}{{\'u}}1
      {Á}{{\'A}}1 {É}{{\'E}}1 {Í}{{\'I}}1 {Ó}{{\'O}}1 {Ú}{{\'U}}1
      {ñ}{{\~n}}1 {Ñ}{{\~N}}1
  },
  title={Salida \thetcbcounter: #2},
  #1
}

%% ---- DATOS DE LA MATERIA ------------------------

\materia{Modelos Estadísticos para el Aprendizaje \\ Automático y Ciencia de Datos}
\actividad{Tema 1.1
Introducción a Conjuntos de Datos Reales y Abiertos \\ Actividad 2 - Proyecto Constructor de Dataset Climático}

\profesor{Dr. José Ismael Ojeda Campaña}
%\fecha{15 de octubre de 2026}    % opcional; por defecto es la fecha de hoy

\begin{document}
\portada
\indicegeneral

%% ---- CONTENIDO ---------------------------------

\section{Introducción}

Se debe dar una introducción de lo que se presenta en el documento \parencite{fisher1936}.

\section{Marco conceptual}

Se debe investigar y agregar teoría y conceptos correspondientes a lo que se está trabajando en la actividad.

\section{Descripción de la actividad}

Descripción de la actividad. Se debe describir de forma detallada el proceso que se realizó para completar la actividad. Se deben añadir figuras, diagramas y código fuente.

\textbf{Objetivo:} Extender el código base proporcionado para ejecutar un Análisis Exploratorio de Datos (EDA) completo sobre el dataset Iris, incorporando segmentación por especies y exportación automatizada de resultados.

\subsection*{Requerimientos del Proyecto}

\begin{description}
    \item[Tarea 1: Extensión del Código Base.] Modificar el script inicial para inspeccionar la estructura y distribución general de los datos:
    \begin{enumerate}
        \item Visualizar las primeras 10 filas del conjunto de datos mediante \texttt{df.head(10)}.
        \item Visualizar las últimas 8 filas mediante \texttt{df.tail(8)}.
        \item Obtener la estructura general, tipos de datos y conteo de valores no nulos con \texttt{df.info()}.
        \item Calcular las estadísticas descriptivas globales (media, desviación estándar, cuartiles) con \texttt{df.describe()}.
    \end{enumerate}

    \item[Tarea 2: Análisis por Especies.] Segmentar la información según la variable objetivo para identificar patrones morfológicos:
    \begin{enumerate}
        \item Mostrar la distribución y número de muestras por cada especie.
        \item Calcular los valores promedio de cada característica agrupados por especie.
        \item Generar un resumen estadístico completo agrupado por especie.
    \end{enumerate}

    \item[Tarea 3: Exportación de Resultados.] Automatizar la persistencia de los datos procesados y los hallazgos obtenidos:
    \begin{enumerate}
        \item Exportar el \textit{DataFrame} completo a un archivo en formato CSV.
        \item Almacenar las tablas estadísticas en un archivo de texto plano (\texttt{.txt}).
        \item Generar de forma automática un reporte resumen con los principales hallazgos del análisis.
    \end{enumerate}
\end{description}

\section {Desarrollo}

Se desarrollaron las actividades anteriores en el Código \ref{cod:iris} que puede verse a continuación:

\begin{codigo}[label=cod:iris]{Exploración Avanzada del Dataset Iris}

import pandas as pd
from sklearn.datasets import load_iris

# Cargar el conjunto de datos Iris
data = load_iris()
df = pd.DataFrame(data.data, columns=data.feature_names)
df['species'] = pd.Categorical.from_codes(data.target, data.target_names)

print("\nPrimeras 10 filas del conjunto de datos:")
print(df.head(10))
print ("\n");
\end{codigo}

\section{Resultados}

Se implementaron satisfactoriamente las actividades correspondientes a las Tareas 1, 2 y 3, cuyo código fuente íntegro se presenta en el Código~\ref{cod:iris}. El cuaderno interactivo puede consultarse en Google Colab:

\begin{center}
    \url{https://colab.research.google.com/drive/1WpdbhrKdbbJqNeG-X2yIbUvCkw1x1dhJ?usp=sharing}
\end{center}

En la Tabla~\ref{tab:promedios_iris} se resumen los valores promedio obtenidos por especie (Tarea 2), donde se observa una clara separación morfológica en las dimensiones del pétalo entre \textit{Iris setosa} y las demás especies.

\begin{table}[htbp]
\centering
\caption{Promedios de las características morfológicas (en cm) agrupados por especie.}
\label{tab:promedios_iris}
\begin{tabular}{lcccc}
\hline
\textbf{Especie} & \textbf{Longitud Sépalo} & \textbf{Ancho Sépalo} & \textbf{Longitud Pétalo} & \textbf{Ancho Pétalo} \\
\hline
\textit{Iris setosa}     & 5.006 & 3.428 & 1.462 & 0.246 \\
\textit{Iris versicolor} & 5.936 & 2.770 & 4.260 & 1.326 \\
\textit{Iris virginica}  & 6.588 & 2.974 & 5.552 & 2.026 \\
\hline
\end{tabular}
\end{table}

Asimismo, la salida completa generada en consola durante la ejecución del script (inspección de registros, tipos de datos y estadísticas descriptivas) se muestra en la Salida~\ref{out:iris}.

\begin{consola}[label=out:iris]{Ejecución en terminal de \texttt{actividad001.py}}
Primeras 10 filas del conjunto de datos:
   sepal length (cm)  sepal width (cm)  petal length (cm)  petal width (cm) species
0                5.1               3.5                1.4               0.2  setosa
1                4.9               3.0                1.4               0.2  setosa
2                4.7               3.2                1.3               0.2  setosa
3                4.6               3.1                1.5               0.2  setosa
4                5.0               3.6                1.4               0.2  setosa
5                5.4               3.9                1.7               0.4  setosa
6                4.6               3.4                1.4               0.3  setosa
7                5.0               3.4                1.5               0.2  setosa
8                4.4               2.9                1.4               0.2  setosa
9                4.9               3.1                1.5               0.1  setosa

\end{consola}

\section{Conclusiones}

Se deben escribir las conclusiones de la actividad

\referencias

\end{document}